% \iffalse meta-comment
% File: sbornik-theorems.dtx (part of the sbornik bundle)
% \fi
%
% \iffalse
%<*driver>
\documentclass[full]{l3doc}
\begin{document}\DocInput{\jobname.dtx}
\end{document}
%</driver>
% \fi
%
% \title{Module \pkg{sbornik-theorems}}
% \author{}\date{}\maketitle
%
% \begin{documentation}
% Theorem module on top of \pkg{amsthm}. Названия переключаются
% через \pkg{babel} (\texttt{\string\addto\string\captionsrussian}).
% \end{documentation}
%
% \begin{implementation}
%    \begin{macrocode}
%<*module>
\bool_new:N \g_sbornik_theorems_bool
\bool_gset_true:N \g_sbornik_theorems_bool

\providecommand\theoremname{Theorem}
\providecommand\lemmaname{Lemma}
\providecommand\corollaryname{Corollary}
\providecommand\definitionname{Definition}
\providecommand\examplename{Example}

\cs_new_protected:Npn \sbornik_theorems_apply:
{
  \bool_if:NT \g_sbornik_theorems_bool
  {
    \RequirePackage{amsthm}
    \theoremstyle{plain}
    \newtheorem{theorem}{\theoremname}[section]
    \newtheorem{lemma}[theorem]{\lemmaname}
    \newtheorem{corollary}[theorem]{\corollaryname}
    \theoremstyle{definition}
    \newtheorem{definition}[theorem]{\definitionname}
    \newtheorem{example}[theorem]{\examplename}
    \addto\captionsrussian{%
      \renewcommand{\theoremname}{Теорема}%
      \renewcommand{\lemmaname}{Лемма}%
      \renewcommand{\corollaryname}{Следствие}%
      \renewcommand{\definitionname}{Определение}%
      \renewcommand{\examplename}{Пример}%
    }
  }
}
%</module>
%    \end{macrocode}
% \end{implementation}
% \Finale
\endinput
